% Options for packages loaded elsewhere
\PassOptionsToPackage{unicode}{hyperref}
\PassOptionsToPackage{hyphens}{url}
\PassOptionsToPackage{dvipsnames,svgnames,x11names}{xcolor}
\documentclass[
  11pt,
  a4paper,
]{article}
\usepackage{xcolor}
\usepackage[margin=2.6cm]{geometry}
\usepackage{amsmath,amssymb}
\setcounter{secnumdepth}{5}
\usepackage{iftex}
\ifPDFTeX
  \usepackage[T1]{fontenc}
  \usepackage[utf8]{inputenc}
  \usepackage{textcomp} % provide euro and other symbols
\else % if luatex or xetex
  \usepackage{unicode-math} % this also loads fontspec
  \defaultfontfeatures{Scale=MatchLowercase}
  \defaultfontfeatures[\rmfamily]{Ligatures=TeX,Scale=1}
\fi
\usepackage{lmodern}
\ifPDFTeX\else
  % xetex/luatex font selection
  \setmainfont[]{Libertinus Serif}
  \setmonofont[Scale=0.85]{Latin Modern Mono}
  \setmathfont[]{Latin Modern Math}
\fi
% Use upquote if available, for straight quotes in verbatim environments
\IfFileExists{upquote.sty}{\usepackage{upquote}}{}
\IfFileExists{microtype.sty}{% use microtype if available
  \usepackage[]{microtype}
  \UseMicrotypeSet[protrusion]{basicmath} % disable protrusion for tt fonts
}{}
\makeatletter
\@ifundefined{KOMAClassName}{% if non-KOMA class
  \IfFileExists{parskip.sty}{%
    \usepackage{parskip}
  }{% else
    \setlength{\parindent}{0pt}
    \setlength{\parskip}{6pt plus 2pt minus 1pt}}
}{% if KOMA class
  \KOMAoptions{parskip=half}}
\makeatother
\usepackage{longtable,booktabs,array}
\newcounter{none} % for unnumbered tables
\usepackage{calc} % for calculating minipage widths
% Correct order of tables after \paragraph or \subparagraph
\usepackage{etoolbox}
\makeatletter
\patchcmd\longtable{\par}{\if@noskipsec\mbox{}\fi\par}{}{}
\makeatother
% Allow footnotes in longtable head/foot
\IfFileExists{footnotehyper.sty}{\usepackage{footnotehyper}}{\usepackage{footnote}}
\makesavenoteenv{longtable}
\setlength{\emergencystretch}{3em} % prevent overfull lines
\providecommand{\tightlist}{%
  \setlength{\itemsep}{0pt}\setlength{\parskip}{0pt}}

\providecommand{\E}{\mathbb{E}}
\providecommand{\N}{\mathbb{N}}
\providecommand{\Z}{\mathbb{Z}}
\providecommand{\R}{\mathbb{R}}
\providecommand{\eps}{\varepsilon}
\providecommand{\ceil}[1]{\left\lceil #1\right\rceil}
\providecommand{\floor}[1]{\left\lfloor #1\right\rfloor}
\AtBeginDocument{\let\setminus\smallsetminus}
\usepackage[most]{tcolorbox}
\newtcolorbox{warning}{colback=red!5,colframe=red!65!black,boxrule=0.8pt,arc=2pt,left=6pt,right=6pt,top=5pt,bottom=5pt}
\usepackage{etoolbox}
\AtBeginEnvironment{longtable}{\small}
\setlength{\emergencystretch}{3em}
\usepackage{fancyhdr}
\pagestyle{fancy}
\fancyhf{}
\fancyhead[L]{\small\itshape Candidate result 2105.15195\_\_00 --- GPT-6 Astra, awaiting review}
\fancyhead[R]{\small\thepage}
\renewcommand{\headrulewidth}{0.2pt}
\usepackage{bookmark}
\IfFileExists{xurl.sty}{\usepackage{xurl}}{} % add URL line breaks if available
\urlstyle{same}
\hypersetup{
  pdftitle={A candidate proof of Conjecture on tightness of the upper bound for all r ≥ 3 in The upper logarithmic density of monochromatic subset sums},
  pdfauthor={github.com/graph-theory-AI --- machine-generated candidate writeup},
  colorlinks=true,
  linkcolor={blue!50!black},
  filecolor={Maroon},
  citecolor={Blue},
  urlcolor={blue!50!black},
  pdfcreator={LaTeX via pandoc}}

\title{A candidate proof of Conjecture on tightness of the upper bound
for all r ≥ 3 in The upper logarithmic density of monochromatic subset
sums}
\author{\href{https://github.com/graph-theory-AI}{github.com/graph-theory-AI}
--- machine-generated candidate writeup}
\date{17 September 2026}

\begin{document}
\maketitle
\begin{abstract}
A prime-based reduction closes the interlacing gap and proves c\_r =
b\_0\^{}(r-1)/(2r) for every r \textgreater= 2. This candidate proof was
generated by GPT-6 Astra in a single model pass for the Graph Theory AI
project. It was selected from the campaign because the model marked it
\texttt{would\_publish:\ true}; it has not been checked by the
adversarial referee, a human mathematician, or a novelty search.
\end{abstract}

\section{The problem and the claimed
result}\label{the-problem-and-the-claimed-result}

This document concerns
\href{https://graph-theory-ai.github.io/graph-conjectures/arxiv/2105.15195__00/}{Conjecture
on tightness of the upper bound for all r ≥ 3} from \emph{The upper
logarithmic density of monochromatic subset sums}.

\subsection{Problem statement}\label{problem-statement}

\emph{The campaign artifact does not contain a separately extracted
statement; see the catalog entry linked above.}

\subsection{Claimed result}\label{claimed-result}

\textbf{Model verdict: proved} (confidence: high).

A prime-based reduction closes the interlacing gap and proves c\_r =
b\_0\^{}(r-1)/(2r) for every r \textgreater= 2.

\emph{The model's own caveat:} The intended denominator is 2r b\_0-r;
the literal 2\^{}r b\_0-r transcription is false for r \textgreater= 3.

\section{Model writeup}\label{model-writeup}

The text below is the model output verbatim, apart from moving its
machine-readable verdict block into the summary above and shifting
Markdown heading levels for this document. Mechanical TeX defects and
equation tags were normalized where needed for compilation.

\subsection{Exact value of the upper logarithmic density
constant}\label{exact-value-of-the-upper-logarithmic-density-constant}

For \(A\subseteq\mathbb N\), write \[
\Sigma(A)=
\left\{\sum_{a\in F}a:
\varnothing\ne F\subseteq A,\ F\text{ finite}\right\},
\] and put \[
\overline d_{\log}(S)
=\limsup_{x\to\infty}
\frac{1}{\log x}\sum_{\substack{n\le x\\n\in S}}\frac1n.
\] Thus the constant in question is \[
c_r=\inf_{\mathbb N=A_1\sqcup\cdots\sqcup A_r}
\max_{i\in[r]}\overline d_{\log}(\Sigma(A_i)).
\]

The key addition to the previous attempt is an additive lemma for
\textbf{dense sets of primes}. It provides the required intervals of
subset sums without requiring any monochromatic interval of integers. I
also give a complete proof of the geometric lemma, including its
boundary accounting.

\subsubsection{Theorem}\label{theorem}

Let \(r\ge2\), and let \(b>1\) be the unique root of \[
b^r-2rb+r-1=0.
\] Then \[
\boxed{\displaystyle
c_r=\frac{b^{r-1}}{2r}
=\frac{1-\frac1{2b}}{1-b^{-r}}
=\left(1-\frac1{2b}\right)
 \left(1+\frac1{2rb-r}\right).}
\qquad\text{(1)}
\]

In particular, this proves the intended tightness conjecture. The
distinction between \(2r\) and \(2^r\) is substantive; it is addressed
explicitly at the end.

\begin{center}\rule{0.5\linewidth}{0.5pt}\end{center}

\subsubsection{1. A geometric covering
lemma}\label{a-geometric-covering-lemma}

For \(r\ge2\) and \(\lambda>1\), define \[
D_{r,\lambda}(q)
=\frac{1-\frac1{\lambda q}}{1-q^{-r}},
\qquad q>1,
\] and \[
C_{r,\lambda}=\inf_{q>1}D_{r,\lambda}(q).
\qquad\text{(2)}
\] Notice that \(0<C_{r,\lambda}<1\). Positivity follows from
\(D_{r,\lambda}(q)\ge1-1/\lambda\), while \(D_{r,\lambda}(q)<1\)
whenever \(q^{r-1}>\lambda\).

\paragraph{Lemma 1}\label{lemma-1}

Suppose \[
0<t_0<t_1<\cdots,\qquad t_k\longrightarrow\infty,
\] and assign each \(k\ge0\) a label \(\gamma_k\in[r]\). Set \[
J_k=[t_k,\lambda t_{k+1}],
\qquad
U_i=\bigcup_{\gamma_k=i}J_k.
\] Then \[
\max_{i\in[r]}
\limsup_{T\to\infty}\frac{|U_i\cap[0,T]|}{T}
\ge C_{r,\lambda},
\qquad\text{(3)}
\] where \(|\cdot|\) is Lebesgue measure.

\paragraph{Proof}\label{proof}

Suppose otherwise. Choose \(D\) satisfying \[
\max_i\limsup_{T\to\infty}\frac{|U_i\cap[0,T]|}{T}
<D<C_{r,\lambda}.
\qquad\text{(4)}
\] For all sufficiently large \(T\), therefore, \[
|U_i\cap[0,T]|\le DT
\qquad(i\in[r]).
\qquad\text{(5)}
\]

No \(U_i\) has an unbounded connected component, since that would give
density \(1\). Every component is consequently a finite interval of the
form \[
C=[t_f,\lambda t_{\ell+1}],
\] where \(f\) and \(\ell\) are the first and last indices of intervals
belonging to that component.

For a component \(C\) of \(U_i\), define \[
e_C=\lambda t_{\ell+1},\qquad
Q_C=\frac{t_{\ell+1}}{t_f},\qquad
a_C=\frac{|C|}{e_C}=1-\frac1{\lambda Q_C},
\] and \[
d_C=\frac{|U_i\cap[0,e_C]|}{e_C}.
\] If \(C^-\) is its preceding component of the same color and \[
R_C=\frac{e_C}{e_{C^-}},
\] then \[
d_C=a_C+\frac{d_{C^-}}{R_C}.
\qquad\text{(6)}
\]

Discard finitely many components so that all components under
consideration have a predecessor and satisfy \(d_C\le D\). For these
components, \[
1-\frac1\lambda<a_C<d_C\le D,
\] and thus \[
1<Q_C<Q_{\max}:=\frac1{\lambda(1-D)}.
\qquad\text{(7)}
\]

Equation (6) gives \[
\log\frac{d_{C^-}}{d_C}
=\log R_C+\log\left(1-\frac{a_C}{d_C}\right)
\le \log R_C+\log\left(1-\frac{a_C}{D}\right).
\qquad\text{(8)}
\]

Consider, on \(1\le q<Q_{\max}\), \[
F(q)=r\log q+
\log\left(1-\frac{1-\frac1{\lambda q}}{D}\right).
\] For \(q>1\), the inequality \(F(q)\ge0\) is equivalent to \[
D\ge D_{r,\lambda}(q),
\] contrary to \(D<C_{r,\lambda}\). Moreover, \[
F(1)<0,\qquad
F(q)\longrightarrow-\infty
\quad\text{as }q\uparrow Q_{\max}.
\] Continuity therefore supplies \(\eta>0\) such that \[
\log\left(1-\frac{a_C}{D}\right)
\le-r\log Q_C-\eta
\qquad\text{(9)}
\] for every component under consideration.

Sum (8) and (9) over the retained components whose last index \(\ell\)
is at most \(N\). Let their number be \(m_N\).

For each color, the left side telescopes. Since \[
1-\frac1\lambda\le d_C\le D
\] for retained components, its total is bounded independently of \(N\):
\[
\sum_C\log\frac{d_{C^-}}{d_C}=O(1).
\qquad\text{(10)}
\] Likewise, telescoping the endpoint ratios gives \[
\sum_C\log R_C\le r\log t_{N+1}+O(1).
\qquad\text{(11)}
\]

The essential covering estimate is \[
\sum_C\log Q_C\ge\log t_{N+1}-O(1).
\qquad\text{(12)}
\] Here is the boundary justification. We have \[
\log Q_C
=\sum_{k=f(C)}^{\ell(C)}
\log\frac{t_{k+1}}{t_k}.
\] Every index \(k\) belongs to the span of its own color-component. The
finitely many discarded components account for only a fixed initial
contribution. At the right boundary \(N\), at most one component of each
color has \(f(C)\le N<\ell(C)\). By (7), the total logarithmic span of
each such component is at most \(\log Q_{\max}\). Thus the total omitted
contribution is bounded independently of \(N\), proving (12).

Combining (8)--(12) yields \[
O(1)
\le r\log t_{N+1}-r\sum_C\log Q_C-\eta m_N+O(1)
\le O(1)-\eta m_N.
\] But \(m_N\to\infty\): there are infinitely many intervals, each
finite component contains only finitely many of them, and no component
is unbounded. This is a contradiction. \(\square\)

\paragraph{\texorpdfstring{Evaluation of
\(C_{r,\lambda}\)}{Evaluation of C\_\{r,\textbackslash lambda\}}}\label{evaluation-of-c_rlambda}

Differentiation gives \[
D_{r,\lambda}'(q)=0
\quad\Longleftrightarrow\quad
q^r-\lambda rq+r-1=0.
\qquad\text{(13)}
\] The polynomial on the right decreases until \(q=\lambda^{1/(r-1)}\),
then increases, and is negative at \(q=1\). It therefore has a unique
root \(b_{r,\lambda}>1\), satisfying \[
b_{r,\lambda}>\lambda^{1/(r-1)}.
\] This root is the unique minimizer in (2). Since \[
b_{r,\lambda}^r-1
=r(\lambda b_{r,\lambda}-1),
\] we obtain \[
C_{r,\lambda}
=\frac{b_{r,\lambda}^{r-1}}{\lambda r}.
\qquad\text{(14)}
\] In particular, \[
\lim_{\lambda\uparrow2}C_{r,\lambda}
=C_{r,2}
=\frac{b^{r-1}}{2r}.
\qquad\text{(15)}
\]

\begin{center}\rule{0.5\linewidth}{0.5pt}\end{center}

\subsubsection{2. A restricted-sum fact modulo a
prime}\label{a-restricted-sum-fact-modulo-a-prime}

The following standard restricted-sum fact is included with a proof so
that the additive reduction is explicit.

\paragraph{Lemma 2}\label{lemma-2}

Let \(p\) be prime and \(B\subseteq\mathbb F_p\), with \(|B|=s\). If \[
1\le k\le s,\qquad k(s-k)\ge p-1,
\] then every element of \(\mathbb F_p\) is a sum of \(k\) distinct
members of \(B\).

\paragraph{Proof}\label{proof-1}

Suppose \(c\in\mathbb F_p\) is not such a sum. The polynomial \[
P(x_1,\ldots,x_k)
=
\left((x_1+\cdots+x_k-c)^{p-1}-1\right)
\prod_{i<j}(x_j-x_i)
\] vanishes on \(B^k\). Indeed, repeated coordinates annihilate the
Vandermonde factor; otherwise the sum differs from \(c\), and Fermat's
theorem annihilates the first factor.

Choose distinct integers \[
0\le d_1<\cdots<d_k\le s-1
\] with \[
\sum_{i=1}^k d_i=p-1+\binom{k}{2}.
\qquad\text{(16)}
\] Such a choice exists because the sums of \(k\) distinct integers from
\(\{0,\ldots,s-1\}\) fill the interval \[
\left[\binom{k}{2},\,
k(s-k)+\binom{k}{2}\right].
\] For example, writing \(p-1=kq+u\), \(0\le u<k\), one may take
\(d_i=i-1+q\) for the first \(k-u\) indices and \(d_i=i+q\) for the last
\(u\).

The coefficient of \(\prod_i x_i^{d_i}\) in \(P\) is \[
\frac{(p-1)!}{\prod_i d_i!}
\prod_{i<j}(d_j-d_i),
\qquad\text{(17)}
\] as follows by expanding the Vandermonde determinant. This is nonzero
in \(\mathbb F_p\), since all \(d_i<p\) and the \(d_i\) are distinct.

On the other hand, a polynomial of total degree \(\sum_i d_i\) with this
coefficient nonzero cannot vanish on \(B^k\), because \(|B|>d_i\) for
every \(i\). This coefficient form of multivariate interpolation follows
by choosing sets \(X_i\subseteq B\) of size \(d_i+1\): the coefficient
equals \[
\sum_{x_i\in X_i}
\frac{P(x_1,\ldots,x_k)}
{\displaystyle\prod_i
 \prod_{y\in X_i\setminus\{x_i\}}(x_i-y)}.
\] It would be zero if \(P\) vanished on \(B^k\), contradicting (17).
\(\square\)

\begin{center}\rule{0.5\linewidth}{0.5pt}\end{center}

\subsubsection{3. Dense prime blocks have long intervals of subset
sums}\label{dense-prime-blocks-have-long-intervals-of-subset-sums}

This is the step that removes the restriction to interval-block
colorings.

\paragraph{Lemma 3 --- Prime-block interval
lemma}\label{lemma-3-prime-block-interval-lemma}

For every fixed \(\kappa>0\), the following holds for all sufficiently
large integers \(n\).

If \(A\) is a set of primes in \([n,2n]\) and \[
|A|\ge \kappa\frac{n}{\log n},
\] then \[
\left[n(\log n)^3,\,
\frac{n^2}{(\log n)^3}\right]\cap\mathbb Z
\subseteq\Sigma(A).
\qquad\text{(18)}
\]

\paragraph{Proof}\label{proof-2}

Write \(m=|A|\). Choose, by Bertrand's postulate, a prime \[
2n<p<4n.
\] All members of \(A\) are distinct modulo \(p\).

Partition \(A\) into a group-building set \(A^{(0)}\), of size
\(\lfloor m/2\rfloor\), and a reserve set \(A^{(1)}\). Put \[
h=\left\lceil\frac{8p}{m}\right\rceil
=O_\kappa(\log n).
\] For sufficiently large \(n\), we have \(h\le m/8\).

\paragraph{Step 1: Build a long arithmetic progression of subset
sums}\label{step-1-build-a-long-arithmetic-progression-of-subset-sums}

As long as at least \(m/4\) elements of \(A^{(0)}\) remain, Lemma 2
supplies an \(h\)-element subset whose sum is \(0\pmod p\). Indeed, if
the remaining size is \(v\), then \[
h(v-h)\ge h\frac m8\ge p.
\] Remove such a subset and continue.

This produces at least \(m/(8h)\) disjoint groups. Each group sum has
the form \[
\ell p,\qquad 1\le\ell\le h,
\] because every group contains \(h\) integers between \(n\) and
\(2n<p\). Consequently, for some \(\ell\), at least \[
L\ge \frac{m}{8h^2}
\] of the groups have the same sum \[
d=\ell p.
\] Their subset sums contain \[
0,d,2d,\ldots,Ld.
\qquad\text{(19)}
\] Here and below \(0\) is allowed temporarily as an empty sum.

Since \(d\ge hn\), the upper endpoint satisfies \[
T:=Ld\ge\frac{mn}{8h}
\gg_\kappa \frac{n^2}{(\log n)^2}.
\qquad\text{(20)}
\]

\paragraph{\texorpdfstring{Step 2: Represent every residue modulo \(d\)
by a small reserve
sum}{Step 2: Represent every residue modulo d by a small reserve sum}}\label{step-2-represent-every-residue-modulo-d-by-a-small-reserve-sum}

All primes in \(A^{(1)}\) exceed \(\ell\), for sufficiently large \(n\).
Thus every occupied residue class modulo \(\ell\) is coprime to
\(\ell\). Choose a residue \(c\pmod\ell\) for which \[
B:=\{a\in A^{(1)}:a\equiv c\pmod\ell\}
\] has size \[
s\ge\frac{m}{2\ell}
\ge\frac{m}{2h}
\gg_\kappa\frac{n}{(\log n)^2}.
\qquad\text{(21)}
\] When \(\ell=1\), the same statements have their usual trivial
interpretation.

Set \[
H=\left\lceil\frac{2p}{s}\right\rceil
=O_\kappa((\log n)^2).
\] For sufficiently large \(n\), \[
H+\ell\le \frac s2.
\] For every \[
t\in\{H,H+1,\ldots,H+\ell-1\},
\] we therefore have \[
t(s-t)\ge H\frac s2\ge p.
\] Lemma 2 says that the sums of \(t\) distinct members of \(B\) cover
all residues modulo \(p\).

Every such sum is also congruent to \(tc\pmod\ell\). Since
\(\gcd(c,\ell)=1\), these \(\ell\) consecutive choices of \(t\) cover
all residues modulo \(\ell\). Also \(\gcd(p,\ell)=1\), because
\(\ell\le h<p\). The Chinese remainder theorem now shows that
\textbf{every residue modulo} \(d=\ell p\) has a representative in
\(\Sigma(B)\) of size at most \[
M:=2n(H+\ell)=O_\kappa(n(\log n)^2).
\qquad\text{(22)}
\]

\paragraph{Step 3: Fill an integer
interval}\label{step-3-fill-an-integer-interval}

Let \(x\) be any integer with \(M\le x\le T\). Choose a reserve sum
\(b_x\le M\) with \[
b_x\equiv x\pmod d.
\] Then \[
x-b_x=qd
\qquad\text{for some }0\le q\le L.
\] Use \(q\) of the disjoint groups from Step 1, together with the
reserve subset representing \(b_x\). These use disjoint elements of
\(A\), so \[
[M,T]\cap\mathbb Z\subseteq\Sigma(A).
\]

Finally, (20) and (22) imply, for sufficiently large \(n\), \[
M\le n(\log n)^3,\qquad
T\ge\frac{n^2}{(\log n)^3}.
\] This proves (18). \(\square\)

The role of primes is precise: they make the reserve elements units
modulo the small integer \(\ell\). No claim about arbitrary dense
integer sets is needed.

\begin{center}\rule{0.5\linewidth}{0.5pt}\end{center}

\subsubsection{4. Applying the geometric lemma to an arbitrary
coloring}\label{applying-the-geometric-lemma-to-an-arbitrary-coloring}

Let \[
\mathbb N=A_1\sqcup\cdots\sqcup A_r
\] be arbitrary. Put \(n_k=2^k\).

By the prime number theorem, for all sufficiently large \(k\), the
interval \([n_k,2n_k]\) contains at least \[
\frac{n_k}{2\log n_k}
\] primes. Some color \(\gamma_k\) therefore contains at least \[
\frac{n_k}{2r\log n_k}
\] of these primes. Apply Lemma 3 with \(\kappa=1/(2r)\). We obtain \[
\left[n_k(\log n_k)^3,\,
\frac{n_k^2}{(\log n_k)^3}\right]\cap\mathbb Z
\subseteq\Sigma(A_{\gamma_k}).
\qquad\text{(23)}
\]

Fix \(1<\lambda<2\), and define \[
t_k=\log n_k+3\log\log n_k.
\] For sufficiently large \(k\), \[
\lambda t_{k+1}
\le 2\log n_k-3\log\log n_k.
\qquad\text{(24)}
\] Indeed, the difference has leading term \((2-\lambda)k\log2\), while
all remaining terms are \(O_\lambda(\log k)\).

It follows from (23)--(24) that every integer whose logarithm lies in \[
J_k=[t_k,\lambda t_{k+1}]
\] belongs to \(\Sigma(A_{\gamma_k})\). Set \[
U_i=\bigcup_{\gamma_k=i}J_k,
\] starting at a sufficiently large index.

We now make the comparison with logarithmic density precise. For
\(E_i=\Sigma(A_i)\), let \[
W_i=\bigcup_{n\in E_i}[\log n,\log(n+1)).
\] For a block \(J_k\) assigned to color \(i\), all its points except
possibly an initial interval of length at most \(e^{-t_k}\) belong to
\(W_i\). Hence \[
|U_i\setminus W_i|
\le\sum_k e^{-t_k}<\infty.
\] For every \(T>0\), consequently, \begin{equation*}
\begin{aligned}
|U_i\cap[0,T]|
&\le |W_i\cap[0,T]|+O(1)\\
&\le\sum_{\substack{n\le e^T\\n\in E_i}}
       \log\left(1+\frac1n\right)+O(1)\\
&\le\sum_{\substack{n\le e^T\\n\in E_i}}\frac1n+O(1).
\end{aligned}
\qquad\text{(25)}
\end{equation*} Thus \[
\overline d_{\log}(\Sigma(A_i))
\ge
\limsup_{T\to\infty}\frac{|U_i\cap[0,T]|}{T}.
\]

Lemma 1 now gives \[
\max_i\overline d_{\log}(\Sigma(A_i))
\ge C_{r,\lambda}.
\] This holds for every \(\lambda<2\). Letting \(\lambda\uparrow2\) and
using (15), \[
\max_i\overline d_{\log}(\Sigma(A_i))
\ge\frac{b^{r-1}}{2r}.
\qquad\text{(26)}
\] The coloring was arbitrary, so \[
c_r\ge\frac{b^{r-1}}{2r}.
\qquad\text{(27)}
\]

This is the unrestricted lower bound; no interval-block assumption
remains.

\begin{center}\rule{0.5\linewidth}{0.5pt}\end{center}

\subsubsection{5. Matching upper bound}\label{matching-upper-bound}

For completeness, the matching construction can also be proved directly.

Let \(b>1\) solve \[
b^r-2rb+r-1=0,
\] and, for all sufficiently large \(k\), put \[
N_k=\left\lceil e^{b^k}\right\rceil.
\] Color \[
[N_k,N_{k+1})\cap\mathbb N
\] by \(k\bmod r\), and color the finite initial interval arbitrarily.

If a monochromatic subset sum has its largest summand in block \(k\),
then it is at least \(N_k\), and it is less than \[
\sum_{n<N_{k+1}}n<N_{k+1}^2.
\] Therefore, apart from finitely many subset sums, \[
\Sigma(A_i)
\subseteq
\bigcup_{k\equiv i\;(\mathrm{mod}\ r)}
[N_k,N_{k+1}^2].
\qquad\text{(28)}
\]

On the logarithmic scale, these intervals have endpoints \[
b^k+O(e^{-b^k}),
\qquad
2b^{k+1}+O(e^{-b^{k+1}}).
\] The errors are summable. Harmonic sums over the corresponding integer
intervals likewise differ from their logarithmic lengths by a summable
error. It is therefore enough to calculate the upper density of \[
V_i=\bigcup_{k\equiv i\;(\mathrm{mod}\ r)}
[b^k,2b^{k+1}].
\]

By the minimizer calculation in Section 1, \[
b>2^{1/(r-1)}.
\] Thus consecutive intervals of the same color are disjoint. At the
right endpoint \(2b^{k+1}\), their accumulated density tends to
\begin{equation*}
\begin{aligned}
\frac{\sum_{j\ge0}(2b-1)b^{k-rj}}{2b^{k+1}}
&=\frac{1-\frac1{2b}}{1-b^{-r}}\\
&=\frac{b^{r-1}}{2r}.
\end{aligned}
\qquad\text{(29)}
\end{equation*} Within an occupied interval the density ratio increases,
and within a gap it decreases. Hence (29) is the upper density of
\(V_i\).

Equation (28) consequently gives \[
\max_i\overline d_{\log}(\Sigma(A_i))
\le\frac{b^{r-1}}{2r}.
\qquad\text{(30)}
\] Together with (27), this proves the theorem.

\begin{center}\rule{0.5\linewidth}{0.5pt}\end{center}

\subsubsection{6. Constant correction and
conclusion}\label{constant-correction-and-conclusion}

The defining equation implies \[
b^r-1=2rb-r.
\] Therefore \[
\frac{b^{r-1}}{2r}
=
\left(1-\frac1{2b}\right)
\left(1+\frac1{2rb-r}\right).
\] For \(r=2\), \(b=2+\sqrt3\), giving \[
c_2=\frac{2+\sqrt3}{4},
\] as required.

For \(r\ge3\), however, \(2^r>2r\), so \[
\left(1-\frac1{2b}\right)
\left(1+\frac1{2^rb-r}\right)
<
\left(1-\frac1{2b}\right)
\left(1+\frac1{2rb-r}\right)
=c_r.
\] Thus the catalog formula with \(2^r\), if interpreted literally, is
false even as an upper bound. The theorem proves the intended \(2r\)
version.

The gap in the previous attempt is closed by Lemma 3: at every dyadic
scale, one color contains enough primes to generate an integer interval
with logarithmic endpoints \[
\log n+O(\log\log n)
\quad\text{and}\quad
2\log n-O(\log\log n).
\] That is exactly what the geometric argument needs. Arbitrary
interlacing of the color classes is therefore allowed, and no unproved
reduction to block colorings is used.

\appendix

\section{Provenance and review
status}\label{provenance-and-review-status}

\begin{warning}

\textbf{UNREFEREED MODEL OUTPUT.} This document was selected solely
because GPT-6 Astra labelled its own result
\texttt{would\_publish:\ true} and returned \texttt{proved} or
\texttt{disproved}. It has not passed the adversarial LLM referee used
for the earlier Sol campaign, has not been checked by a human
mathematician, and has not been checked for novelty. Treat every
mathematical and bibliographic claim as unverified.

\end{warning}

{\def\LTcaptype{none} % do not increment counter
\begin{longtable}[]{@{}
  >{\raggedright\arraybackslash}p{(\linewidth - 2\tabcolsep) * \real{0.5000}}
  >{\raggedright\arraybackslash}p{(\linewidth - 2\tabcolsep) * \real{0.5000}}@{}}
\toprule\noalign{}
\endhead
\bottomrule\noalign{}
\endlastfoot
Catalog id & \texttt{2105.15195\_\_00} \\
Catalog entry &
\href{https://graph-theory-ai.github.io/graph-conjectures/arxiv/2105.15195__00/}{Conjecture
on tightness of the upper bound for all r ≥ 3} \\
Source corpus & arxiv \\
Campaign leg & selected second attempts (\texttt{attacks\_retry}) \\
Source paper / entry & The upper logarithmic density of monochromatic
subset sums \\
Model verdict & \textbf{proved} (confidence: high) \\
Selection criterion & \texttt{would\_publish:\ true} and verdict
\texttt{proved} or \texttt{disproved} \\
Model & \texttt{gpt-6-astra}, reasoning effort \texttt{max},
\texttt{mode=pro}, flex service tier \\
Original artifact &
\texttt{attacks\_retry/2105.15195\_\_00/output.md} \\
Independent review & \textbf{None} \\
arXiv & \href{https://arxiv.org/abs/2105.15195}{arXiv:2105.15195} \\
Previous attempt & \texttt{attacks/2105.15195\_\_00}: gpt-5.6-sol
verdict \texttt{partial} \\
Model's caveats & The intended denominator is 2r b\_0-r; the literal
2\^{}r b\_0-r transcription is false for r \textgreater= 3. \\
\end{longtable}
}

The generated Markdown and LaTeX sources for this PDF are stored under
\texttt{to\_review\_astra/src/2105.15195\_\_00/}.

\end{document}
